% Options for packages loaded elsewhere
\PassOptionsToPackage{unicode}{hyperref}
\PassOptionsToPackage{hyphens}{url}
\documentclass[
  ignorenonframetext,
]{beamer}
\newif\ifbibliography
\usepackage{pgfpages}
\setbeamertemplate{caption}[numbered]
\setbeamertemplate{caption label separator}{: }
\setbeamercolor{caption name}{fg=normal text.fg}
\beamertemplatenavigationsymbolsempty
% remove section numbering
\setbeamertemplate{part page}{
  \centering
  \begin{beamercolorbox}[sep=16pt,center]{part title}
    \usebeamerfont{part title}\insertpart\par
  \end{beamercolorbox}
}
\setbeamertemplate{section page}{
  \centering
  \begin{beamercolorbox}[sep=12pt,center]{section title}
    \usebeamerfont{section title}\insertsection\par
  \end{beamercolorbox}
}
\setbeamertemplate{subsection page}{
  \centering
  \begin{beamercolorbox}[sep=8pt,center]{subsection title}
    \usebeamerfont{subsection title}\insertsubsection\par
  \end{beamercolorbox}
}
% Prevent slide breaks in the middle of a paragraph
\widowpenalties 1 10000
\raggedbottom
\AtBeginPart{
  \frame{\partpage}
}
\AtBeginSection{
  \ifbibliography
  \else
    \frame{\sectionpage}
  \fi
}
\AtBeginSubsection{
  \frame{\subsectionpage}
}
\usepackage{iftex}
\ifPDFTeX
  \usepackage[T1]{fontenc}
  \usepackage[utf8]{inputenc}
  \usepackage{textcomp} % provide euro and other symbols
\else % if luatex or xetex
  \usepackage{unicode-math} % this also loads fontspec
  \defaultfontfeatures{Scale=MatchLowercase}
  \defaultfontfeatures[\rmfamily]{Ligatures=TeX,Scale=1}
\fi
\usepackage{lmodern}
\ifPDFTeX\else
  % xetex/luatex font selection
\fi
% Use upquote if available, for straight quotes in verbatim environments
\IfFileExists{upquote.sty}{\usepackage{upquote}}{}
\IfFileExists{microtype.sty}{% use microtype if available
  \usepackage[]{microtype}
  \UseMicrotypeSet[protrusion]{basicmath} % disable protrusion for tt fonts
}{}
\makeatletter
\@ifundefined{KOMAClassName}{% if non-KOMA class
  \IfFileExists{parskip.sty}{%
    \usepackage{parskip}
  }{% else
    \setlength{\parindent}{0pt}
    \setlength{\parskip}{6pt plus 2pt minus 1pt}}
}{% if KOMA class
  \KOMAoptions{parskip=half}}
\makeatother
\setlength{\emergencystretch}{3em} % prevent overfull lines
\providecommand{\tightlist}{%
  \setlength{\itemsep}{0pt}\setlength{\parskip}{0pt}}
\usepackage{bookmark}
\IfFileExists{xurl.sty}{\usepackage{xurl}}{} % add URL line breaks if available
\urlstyle{same}
\hypersetup{
  hidelinks,
  pdfcreator={LaTeX via pandoc}}

\author{\texorpdfstring{}{}}
\date{}

\begin{document}

\begin{frame}[fragile]{Experimental Setup}
\protect\phantomsection\label{experimental-setup}
This section details the complete experimental configuration used to
generate the results in this study. Every parameter value below is
injected programmatically from \texttt{config.yaml} and the analysis
pipeline --- no value is hardcoded in the manuscript source.

\begin{block}{Problem Definition}
\protect\phantomsection\label{problem-definition}
The optimization target is the quadratic objective function:

\begin{equation}
\label{eq:quadratic_objective}
f(x) = \frac{1}{2} x^T A x - b^T x
\end{equation}

with \(A = [[1.0]]\) and \(b = [1.0]\), yielding the analytical optimum
at \(x^* = 1.0\) with \(f(x^*) = -0.5\).
\end{block}

\begin{block}{Parameter Space}
\protect\phantomsection\label{parameter-space}
The experiment systematically varies the gradient descent step size
across 6 values:

\begin{itemize}
\tightlist
\item
  \(\\alpha = 0.01\) (conservative)
\item
  \(\\alpha = 0.1\) (moderate)
\item
  \(\\alpha = 0.5\) (aggressive)
\item
  \(\\alpha = 1.0\) (very aggressive)
\item
  \(\\alpha = 1.5\)
\item
  \(\\alpha = 2.5\)
\end{itemize}

All runs start from the initial point \(x_0 = 0.0\) and use a
convergence tolerance of \(\|{\nabla f}\| < 10^{-8}\) with a maximum
iteration limit of \(N_{\max} = 1000\).
\end{block}

\begin{block}{Numerical Stability Grid}
\protect\phantomsection\label{numerical-stability-grid}
To validate robustness, the optimizer is exercised across a grid of 8
starting points and 6 step sizes, producing 48 total evaluations. This
comprehensive sweep confirms that convergence is not an artifact of a
narrow parameter choice.
\end{block}

\begin{block}{Dimensional Scaling}
\protect\phantomsection\label{dimensional-scaling}
Performance benchmarking spans problem dimensions
\(d \in \{1, 2, 5, 10, 20, 50\}\), from the scalar case (\(d = 1\)) to
moderate dimensionality (\(d = 50\)), using identity-Hessian quadratics
to isolate algorithmic scaling from problem conditioning effects.
\end{block}

\begin{block}{Computational Environment}
\protect\phantomsection\label{computational-environment}
\begin{itemize}
\tightlist
\item
  \textbf{Python}: 3.12.11
\item
  \textbf{NumPy}: 2.4.1
\item
  \textbf{Platform}: Darwin arm64
\item
  \textbf{Generated}: 2026-03-22T21:31:47Z
\end{itemize}
\end{block}
\end{frame}

\end{document}
